%MB9T5
\begin{vidu}%KNTT
Trong một thí nghiệm về giao thoa ánh sáng với hai khe Young, khoảng cách giữa hai khe hẹp là $a=2\ mm$, khoảng cách giữa mặt phẳng chứa hai khe với màn quan sát là $D=1,2\ m$. Khe sáng hẹp phát đồng thời hai bức xạ đơn sắc màu đỏ $\lambda_1=0,6\ \mu m$ và màu lục $\lambda_2=0,5\ \mu m$.
\begin{itemize}
\item[1.] Khoảng vân của ánh sáng màu đỏ là $i_1=..........\ mm$ và màu lục là $i_2=..........\ mm$. 
\item[2.] Khoảng cách ngắn nhất giữa hai vân sáng cùng màu với vân sáng trung tâm là .......... mm
\item[3.] Khoảng cách ngắn nhất giữa hai vân sáng bậc ba cùng phía với vân trung tâm là .......... mm
\item[4.] Số vân sáng trùng nhau trên bề rộng $3\ cm$ của vùng giao thoa là ..........
\item[5.] Biết M, N có tọa độ tương ứng là 2 mm và 8 mm. Số vân sáng trên đoạn MN là ..........
\item [6.] Số vân đơn sắc giữa hai vân sáng trùng nhau  là ..........
\end{itemize}
\loigiai{
\begin{itemize}
\item [1.] Khoảng vân của hai ánh sáng:
$$i_1=\dfrac{\lambda_1.D}{a}=0,36\ mm;\ i_2=\dfrac{\lambda_2.D}{a}=0,3\ mm$$
\item [2.] Khoảng cách ngắn nhất giữa hai vân cùng màu với trung tâm là khoảng cách giữa hai vân sáng trùng: $\dfrac{i_1}{i_2}=\dfrac{6}{5}\Rightarrow i_{\equiv}=5i_1=1,8\ mm$
\item [3.] Khoảng cách ngắn nhất giữa hai vân bậc 3 cùng phía: $\Delta x=3|i_1-i_2|=0,18\ mm$ 
\item [4.] Số vân sáng trùng nhau trên bề rộng 3 mm: $N_{\equiv}=2.\left[\dfrac{\ell}{2i_{\equiv}}\right]+1=17$
\item [5.] Số vân sáng trên MN: 
\begin{itemize}
\item Số vân sáng của $\lambda_1$: $\dfrac{x_M}{i_1}\le k_1\le\dfrac{x_M}{i_1}\Rightarrow 5,5\le k_1\le 22,2\Rightarrow N_1=17$
\item Số vân sáng của $\lambda_2$: $\dfrac{x_M}{i_2}\le k_2\le\dfrac{x_M}{i_2}\Rightarrow 6,6\le k_2\le 26,6\Rightarrow N_2=20$
\item Số vân sáng trùng: $\dfrac{x_M}{i_{\equiv}}\le k_{\equiv}\le\dfrac{x_M}{i_{\equiv}}\Rightarrow 1,1\le k_{\equiv}\le 4,4\Rightarrow N_{\equiv}=3$
\item Số vân sáng: $N_s=N_1+N_2-N_{\equiv}=34$
\end{itemize}
\item [6.] Ta có: $i_{\equiv}=5i_1=6i_2$ nên giữa hai vân sáng trùng có 4 vân đơn sắc của $\lambda_1$ và 5 vân đơn sắc của $\lambda_2$. Tổng cộng có 9 vân đơn sắc.

\end{itemize}
}
\end{vidu} \haidong{20}

\loai{Cho bậc của hai vân trùng. Tìm bước sóng}
\begin{vidu}%[3P5-4.1K][Loai3]
Trong thí nghiệm Young, khoảng cách giữa hai khe là 1\ mm, khoảng cách từ hai khe đến màn là 2 m. Chiếu vào hai khe ánh sáng hỗn tạp gồm hai bức xạ có $\lambda_1=0,72\ \mu m$ và $\lambda_2$, người ta thấy vân sáng bậc 3 của bức xạ $\lambda_2$ trùng với vân sáng bậc 2 của bức xạ $\lambda_1$. Bước sóng  $\lambda_2$ có giá trị bằng
\choice
{$0,54\ \mu m$}
{$0,43\ \mu m$}
{\True $0,48\ \mu m$}
{$0,45\ \mu m$}
\loigiai{
\begin{itemize}
\item Ta có: $x=3\dfrac{\lambda_2D}{a}=2\dfrac{\lambda_1D}{a}\Rightarrow \lambda_2=\dfrac{2\lambda_1}{3}=0,48\ \mu m$.
\end{itemize}
}
\end{vidu} \haidong{20}
\loai{Tìm bậc của vân đơn sắc tại vị trí trùng}
\begin{vidu}%[3P5-4.1K][Loai4]
Ánh sáng được dùng trong thí nghiệm giao thoa gồm hai ánh sáng đơn sắc ánh sáng lục có bước sóng $\lambda_1=0,50\ \mu m $ và ánh sáng đỏ có bước sóng $\lambda_2=0,75\ \mu m $. Vân sáng trùng gần vân trung tâm nhất ứng với vân sáng đỏ bậc
\choice
{5}
{6}
{4}
{\True 2}
\loigiai{
\begin{itemize}
\item Ta có 2 vân trùng nhau khi: ${k_1}{i_1}={k_2}{i_2}\Rightarrow {k_1}\lambda_1={k_2}\lambda_2\Rightarrow \dfrac{k_1}{k_2}=\dfrac{\lambda_2}{\lambda_1}=\dfrac{3}{2} $.
\item Để là vân trùng gần vân trung tâm nhất thì k phải nhỏ nhất $\Rightarrow {k_2}=2 $.
\end{itemize}
}
\end{vidu} \haidong{20}
\loai{Tìm số vân đơn sắc giữa hai vân trùng}
\begin{vidu}%[3P5-4.1K][Loai5]
Trong thí nghiệm Young về giao thoa ánh sáng, nguồn sáng phát đồng thời hai ánh sáng đơn sắc $\lambda_1,\ \lambda_2 $ có bước sóng lần lượt là 0,48 mm và 0,60 mm. Trên màn quan sát, trong khoảng giữa hai vân sáng gần nhau nhất và cùng màu với vân sáng trung tâm có
\choice
{\True 4 vân sáng $\lambda_1 $ và 3 vân sáng $\lambda_2 $}
{5 vân sáng $\lambda_1 $ và 4 vân sáng $\lambda_2 $}
{4 vân sáng $\lambda_1 $ và 5 vân sáng $\lambda_2 $}
{3 vân sáng $\lambda_1 $ và 4 vân sáng $\lambda_2 $}
\loigiai{
\begin{itemize}
\item Ta có 2 vân trùng nhau khi: ${k_1}{i_1}={k_2}{i_2}\Rightarrow {k_1}\lambda_1={k_2}\lambda_2\Rightarrow \dfrac{k_1}{k_2}=\dfrac{\lambda_2}{\lambda_1}=\dfrac{5}{4} $.
\item Để là vân trùng gần vân trung tâm nhất thì k phải nhỏ nhất, khi đó ${k_1}=5;\ {k_2}=4 $.
\item Vậy giữa 2 vân sáng gần nhau nhất cùng màu với vân trung tâm có 4 vân $\lambda_1 $ và 3 vân $\lambda_2 $.
\end{itemize}
}
\end{vidu} \haidong{20}
\loai{Tìm số vân trùng trên trường giao thoa}
\begin{vidu}%[3P5-4.1K][Loai6]
Trong thí nghiệm Young về giao thoa ánh sáng, nguồn sáng S chiếu đến hai khe \bth{S_1,\ S_2} đồng thời hai bức xạ đơn sắc có bước sóng \bth{\lambda_1=0,5\microm} và \bth{\lambda_2=0,75\microm}. Khoảng cách giữa hai khe là 1\ mm, khoảng cách từ hai khe tới màn là 2\ m. Xét trên bề rộng vùng giao thoa \bth{L=32,7\ mm}, số vân sáng trùng nhau của hai bức xạ là
\choice
{\True 11}
{23}
{16}
{27}
\loigiai{
\begin{itemize}
\item Những vị trí mà vân sáng của hai bức xạ trùng nhau sẽ có cùng tọa độ trên màn
\begin{align*}
x_s^{\lambda_1}=x_s^{\lambda_2}\Leftrightarrow k_1.\dfrac{\lambda_1.D}{a}=k_2.\dfrac{\lambda_2.D}{a}\Rightarrow \dfrac{k_1}{k_2}=\dfrac{\lambda_2}{\lambda_1}=\dfrac{0,75}{0,5}=\dfrac{3}{2}.
\end{align*}
\item Vì \bth{k_1,\ k_2\in\mathbb{Z}} nên \bth{k_1=3n;\ k_2=2n} với \bth{n\in \mathbb{Z}}. Chọn 1 bức xạ để tính toán.
\begin{itemize}
\item Ta có:
$x_s=x_s^{\lambda_1}=k_1\dfrac{\lambda_1.D}{a}=3n\dfrac{\lambda_1.D}{a}.$
\item Số vân trùng trên bề rộng vùng giao thoa:
\begin{align*}
-\dfrac{32,7}{2}\leq 3n\dfrac{\lambda_1.D}{a}\leq \dfrac{32,7}{2}\Leftrightarrow -\dfrac{32,7}{2}\leq 3n\dfrac{0,5.2}{1}\leq \dfrac{32,7}{2}\Rightarrow -5,45\leq n\leq 5,45.
\end{align*}
\end{itemize}
\item Vậy có 11 vân sáng trùng nhau. 
\end{itemize}
}
\end{vidu} \haidong{20}
\loai{Tìm số vân sáng trên trường giao thoa}
\begin{vidu}%[3P5-4.1G][Loai7] 
Trong thí nghiệm giao thoa Young, nếu chiếu bức xạ có bước sóng $\lambda _1=0,4\ \mu m$ thì trên bề rộng L người ta thấy 31 vân sáng, nếu thay bước sóng $\lambda _1$ bằng bức xạ có bước sóng $\lambda _2=0,6\ \mu m$ thì người ta thấy có 21 vân sáng. Biết trong cả hai trường hợp thì ở hai điểm ngoài cùng của khoảng L đều là vân sáng. Nếu chiếu đồng thời hai bức xạ thì trên bề rộng L quan sát được
\choice
{\True 41 vân sáng}
{40 vân sáng}
{52 vân sáng}
{36 vân sáng}
\loigiai{
\begin{itemize}
\item Trên bề rộng L có 31 vân sáng của bức xạ $\lambda _1$ nên: $L=\left({31-1}\right).i_1=30.i_1$.
\item Tổng số vân của 2 bức xạ trên bề rộng L: $N=N_1+N_2=31+21=52$ vân.
\item Điều kiện trùng nhau của hai bức xạ: $\dfrac{k_1}{k_2}=\dfrac{i_2}{i_1}=\dfrac{\lambda _2}{\lambda _1}=\dfrac{0,6}{0,4}=\dfrac{3}{2}\Rightarrow k_1=3$.
\item Khoảng vân trùng là: $i_\equiv=\dfrac{k_1\lambda _1D}{a}=k_1.i_1=3.i_1$.
\item Số vân trùng nhau trên bề rộng L: $N_\equiv=1+2.\left[{\dfrac{L}{2i_\equiv}}\right]=1+2.\left[{\dfrac{30i_1}{2.3i_1}}\right]=1+2.\left[5\right]=11$ vân.
\item Số vân sáng quan sát được: $N_S=N-N_\equiv=52-11=41$ vân.
\end{itemize}}
\end{vidu} \haidong{20}
\loai{Tìm số vân trùng trên đoạn thẳng}
\begin{vidu}%[3P5-4.1K][Loai8]
\nguon{QG - 2017} Trong thí nghiệm Young về giao thoa ánh sáng, hai khe được chiếu bằng ánh sáng gồm hai thành phần đơn sắc có bước sóng $\lambda =0,6\ \mu m$ và $\lambda'=0,4\ \mu m$. Trên màn quan sát, trong khoảng giữa hai vân sáng bậc 7 của bức xạ có bước sóng $\lambda$, số vị trí cho vân trùng nhau giữa hai bức xạ là
\choice
{8}
{5}
{6}
{\True 7}
\loigiai{
\begin{itemize}
\item Điều kiện để hai hệ vân trùng nhau
$\dfrac{\lambda}{\lambda'}=\dfrac{k'}{k}=\dfrac{6}{4}=\dfrac{3}{2}$.
\item Ta có thể lập bảng sau:
\begin{bang}[tabularx={X|X|X|X|X|X}]{}
k&0&$\pm2$&$\pm4$&$\pm6$&$\pm8$\\\hline
k'&0&$\pm3$&$\pm6$&$\pm9$&$\pm12$
\end{bang}
\item Từ bảng trên ta thấy rằng giữa hai vân sáng bậc 7 của bước sóng $\lambda$ có 7 vị trí cho vân trùng nhau của hai bức xạ (lưu ý rằng vị trí trung tâm được tính là 1).
\end{itemize}
}
\end{vidu} \haidong{20}
\loai{Tìm số vân sáng trên đoạn thẳng}
\begin{vidu}%[3P5-4.1G][Loai9]
Trong thí nghiệm Young về giao thoa ánh sáng, thực hiện đồng thời với hai ánh sáng đơn sắc. Khoảng vân giao thoa trên màn lần lượt là 1,2 mm và 1,8 mm. Trên màn quan sát, gọi M, N là hai điểm ở cùng một phía so với vân trung tâm và cách vân trung tâm lần lượt là 6\ mm và 20\ mm. Trên đoạn MN, quan sát được bao nhiêu vạch sáng?
\choice
{19}
{\True 16}
{20}
{18}
\loigiai{ 
\begin{itemize}
\item Số vị trí vân sáng trùng nhau trên MN:\\
$\left\{\begin{aligned}& x=k_1i_1=k_2i_2=k_1.1,2=k_2.1,8\ mm\Rightarrow \dfrac{k_1}{k_2}=\dfrac{3}{2}\Rightarrow \left\{\begin{aligned}& k_1=3n \\
& k_2=2n 
\end{aligned}\right. \\
& x=3n.1,2\ mm=3,6n\ mm\Rightarrow \left\{\begin{aligned}& 6\le x\le 20\Rightarrow 1,7\le n\le 5,6 \\
& \Rightarrow n=2,\ 3,\ 4,\ 5:\text{ số vị trí trùng: 4} 
\end{aligned}\right. 
\end{aligned}\right.$
\item Số vân sáng của hệ 1 và hệ 2 trên MN lần lượt được xác định như sau:\\
$\left\{\begin{aligned}& 6\le x=k_1i_1=k_1.1,2\le 20\Rightarrow 5\le k_1\le 16,7\Rightarrow k_1=5,...,16:\text{ số giá trị } k_1 \text{ là 12} \\
& 6\le x=k_2i_2=k_2.1,8\le 20\Rightarrow 3,3\le k_2\le 11,1\Rightarrow k_2=4,...,11:\text{ số giá trị } k_2 \text{ là 8} 
\end{aligned}\right.$\\
$\Rightarrow \text{Số vạch sáng}=12+8-4=16$.
\end{itemize}
}
\end{vidu} \haidong{20}
\loai{Tìm số vân đơn sắc}
\begin{vidu}%[3P5-4.1G][Loai10]
Chiếu đồng thời hai ánh sáng đơn sắc có bước sóng $0,54\ \mu m$ và $0,72\ \mu m$ vào hai khe của thí nghiệm Young. Biết khoảng cách giữa hai khe 0,8 mm, khoảng cách từ hai khe tới màn 1,8 m. Trong bề rộng trên màn 2  cm (vân trung tâm ở chính giữa), số vân sáng của hai bức xạ không có màu giống màu của vân trung tâm là
\choice
{\True 20}
{5}
{25}
{30}
\loigiai{ 
\begin{itemize}
\item Ta có: $i_1=\dfrac{\lambda_1D}{a}=1,215\ mm$\quad $i_2=\dfrac{\lambda_2D}{a}=1,62\ mm$
\item Khoảng vân trùng: $\dfrac{i_2}{i_1}=\dfrac{1,62}{1,215}=\dfrac{4}{3}\Rightarrow i_\equiv =4i_1=3i_2=4.1,215=4,86\ mm$
\item Số vân tương ứng:
\begin{align*}
&N_\equiv =2\left[{\dfrac{0,5L}{i_\equiv}}\right]+1=2\left[{\dfrac{0,5.20}{4,86}}\right]+1=5\\
&N_1=2\left[{\dfrac{0,5L}{i_1}}\right]+1=2\left[{\dfrac{0,5.20}{1,215}}\right]+1=17\\
&N_2=2\left[{\dfrac{0,5L}{i_2}}\right]+1=2\left[{\dfrac{0,5.20}{1,62}}\right]+1=13
\end{align*}
\item Số vân sáng khác màu với vân trung tâm $17+13-2.5=20$.
\end{itemize}
}
\end{vidu} \haidong{20}
 
